\section{Check 4: Response}\label{sec:response}

\emph{What a per-flow FPR hides.} A per-flow Benign FPR of a few percent looks modest in a table, but at a gateway that blocks host pairs it means blocked innocent hosts. False alarms already dominate security analysts' work~\cite{alahmadi2022falsepositives}, and a few false positives can have strong consequences once detections drive IP blocking~\cite{garcia2023labeling,akbari2020atmos}. We therefore count false blocks per \emph{host-pair window}, a host pair within one 60\,s test window: a block is false if the pair carries no attack flow in that window, and the pairs without an attack flow are the eligible benign host-pair windows.

\emph{Policies.} With enforcement on, the gateway installs OVS drop rules for both directions of a host pair once the pair has $k$ attack verdicts, optionally ignoring MITM-ARP verdicts, which an ARP watcher covers (below); the harness resets the counts after each window. On validation we meant to choose among ten policies ($k\in\{1,2,3,5,10\}$, MITM-ARP verdicts ignored or not) the one with the fewest false blocks among those within 0.005 of the best share of attack flows that start after their pair is blocked. This rule was degenerate: all ten policies are nearly equal on that share, so it picked the largest $k$ ($k=10$, MITM-ARP ignored; 1 false block on validation, against 128 at $k=1$). We therefore report the trade-off over $k$ (Table~\ref{tab:kcurve}).

\begin{table}[t]
\centering
\caption{Trade-off over $k$, simulated per 60\,s window on the verdicts logged in run~1 (the enforced $k=1$ run, a separate replay, gave 60 false blocks). Ign.: MITM-ARP verdicts ignored. False: false blocks in the ten 60\,s windows. Class columns: attacker pairs blocked, of the number in parentheses (BF BruteForce, DS DDoS-SYN, DH DoS-HTTP, MA MITM-ARP, Mi Mirai, Re Recon); Att.: their sum. Time: median, over the attacker pairs blocked at that $k$, of the time from the pair's first flow to its block.}
\label{tab:kcurve}
\footnotesize
\setlength{\tabcolsep}{2.4pt}
\begin{tabular}{rcrcccccccrr}
\hline
 & & & BF & DS & DH & MA & Mi & Re & Web & Att. & Time \\
$k$ & Ign. & False & (1) & (6) & (1) & (7) & (7) & (1) & (3) & (26) & (s) \\
\hline
1  & no  & 64 & 1 & 6 & 1 & 0 & 7 & 1 & 3 & 19 & 19.0 \\
1  & yes & 37 & 1 & 6 & 1 & 0 & 7 & 1 & 3 & 19 & 19.0 \\
2  & yes &  5 & 1 & 6 & 1 & 0 & 6 & 1 & 3 & 18 & 33.8 \\
3  & yes &  1 & 1 & 6 & 1 & 0 & 1 & 1 & 3 & 13 & 15.8 \\
5  & yes &  0 & 1 & 6 & 1 & 0 & 0 & 1 & 1 & 10 & 15.3 \\
10 & yes &  0 & 1 & 6 & 1 & 0 & 0 & 1 & 0 &  9 & 15.3 \\
20 & yes &  0 & 0 & 6 & 1 & 0 & 0 & 1 & 0 &  8 & 15.3 \\
\hline
\end{tabular}
\end{table}

\emph{The trade-off.} The DDoS-SYN, DoS-HTTP and Recon attackers are blocked at every $k$, so $k$ trades false blocks against the low-rate attackers (Table~\ref{tab:kcurve}): most Mirai bots are lost at $k=3$ and all from $k=5$, the web attackers drop to 1 of 3 at $k=5$ and to none from $k=10$, and the brute-force attacker is lost at $k=20$. Counting MITM-ARP verdicts adds false blocks at small $k$.

\emph{Enforced runs.} We ran both ends of the curve live over the ten windows (blocks cleared per window); no live window contains the .174 flood, so the audit's corrections do not affect these counts. Blocking on the first verdict blocked 19 of the 26 attacker pairs, including all 7 Mirai bots, but also 60 of 1,674 eligible benign host-pair windows (3.6\%), in 9 of the 10 windows. Almost all of these false blocks (56) hit pairs of uninvolved hosts, and 21 were triggered by MITM-ARP verdicts. The $k=10$ policy (MITM-ARP ignored) blocked only the six DDoS attackers, the DoS-HTTP attacker, the brute-force attacker and the port scanner, with no false blocks; it missed the 7 Mirai bots and the 3 web attackers. In the $k=1$ run the DDoS-SYN window was replayed at only 52\% of the original rate, so flow splitting (Section~\ref{sec:fidelity}) may inflate its counts there.

\emph{Persistent counters.} To separate these losses from the per-window reset, we replayed the rule on the deployed tree's verdicts over each capture's whole temporal test period (5.9\,h; 55 attacker and 3,410 benign-only host pairs; audit-corrected labels), with counters reset every 60\,s or never cleared. The Mirai loss belongs to the flow-count rule, not to the classifier (at $k=1$ all 7 bots were blocked live) or the reset: each bot's flood forms only one or two long flows, so even persistent counters block at most one bot from $k=3$. The web loss at $k=5$ belongs to the harness: persistent counters block all 7 web attacker pairs, the reset only 3. The reset also hid false blocks: with persistent counters, $k=3$, 5 and 10 cause 21, 6 and 2 false blocks instead of 2, 0 and 0 (one more each with the original labels). At $k=1$ the replay falsely blocks 3.6\% of the benign-only pairs, about 21 per hour. These rates are not comparable with the live counts, which use 60\,s host-pair windows and attack-dense traffic. A deployed policy therefore needs persistent counters and a criterion that does not count flows alone.

\emph{ARP watcher.} For MITM-ARP, which flow features cannot see (Section~\ref{sec:splits}), a watcher in the style of arpwatch~\cite{arpwatch} alerts when an ARP reply binds the router IP (.1) to a MAC other than the known one. It alerted once, 1.7\,s into the MITM-ARP window, and in none of the other nine: a functional check, not a false-alarm rate.

\begin{figure}[t]
\centering
\includegraphics[width=\columnwidth]{late_blocks.pdf}
\caption{Share of attack flows that start after their host pair is blocked ($k=1$, MITM-ARP verdicts ignored), with each verdict stamped at flow end, as an offline evaluation does, or 15\,s later, when the live gateway exports the flow (idle timeout). Top: pooled over run~1's ten 60\,s windows, counters reset per window (logged live verdict times: 11.1\%). Below: the 5.9\,h test periods with persistent counters and corrected labels, pooled and per class (MITM-ARP and Mirai, 35 flows, omitted).}
\label{fig:lateblocks}
\end{figure}

\emph{Late blocks.} Offline policy scores overstate prevention (Fig.~\ref{fig:lateblocks}). With each verdict stamped at flow end, as an offline evaluation does, about 97\% of the attack flows in the live windows start after their pair is blocked; with the verdict delayed by the 15\,s idle timeout, or at the logged live verdict time, only 11\% do. In these short windows the export delay, not the window cut, causes the drop: in the DDoS-SYN window the first verdict arrives 15\,s in, when 82\% of the attack flows have already started. Over the 5.9\,h test periods the pooled share falls to 20\%, a value dominated by the trimmed DDoS-SYN capture, whose test period lasts only 23\,s. A long test period does not by itself protect a class, though: the loss is small for DoS-HTTP and the web attacks but large for the port scan and brute force. Blocks still cut the later packets of long attack flows (34--71\% of the DDoS-SYN and DoS-HTTP windows' packets).
